\subsection{Dart}

\subsection*{識別子}
Dartの識別子には以下のように大きく３種類に分けることが出来る。\par
\begin{itemize}
  \item \textbf{UpperCamelCase}：名前では、各単語の最初の文字が大文字になり、最初の文字も含まれる。
  \item \textbf{lowerCamelCase}：名前では、頭文字を覗き、各単語の最初の文字を大文字にする
  \item \textbf{snake\_case}：頭字語であっても、小文字のみを使用し、単語は\_で区切る。
\end{itemize}
これを元に型やファイルなどの識別子を定める。

\subsubsection*{・型にはUpperCamelCaseを用いる。}
クラス、列挙型、typedefs、型パラメータは、各単語の最初の文字を大文字にし（最初の単語を含む）、セパレータを使用しない。\par
これには、メタデータ注釈で使用されることを意図したクラスも含む。\par

アノテーションクラスのコンストラクタがパラメータを取らない場合、それ用に別のloweCamelCase定数を作成した方が良い。\par

\subsubsection*{・拡張子には、UpperCamelCaseを使う。}
タイプ同様、拡張子は各単語の最初の文字を大文字に（最初の単語を含む）、セパレータを使用しない。\par





\subsubsection*{・ライブラリ、パッケージ、ディレクトリ、ソースファイルにはsnake\_caseを使う}
ファイルシステムによっては大文字と小文字が区別されないため、
多くのプロジェクトでファイル名を全て小文字にする必要がある。
区切り文字を使用することで、
そのような形式でも名前を読むことができる。
アンダースコア（\_）をセパレータとして使用すると、
名前が有効なDart識別子であることが保証される。これは、
言語が後にシンボリックインポートをサポートする場合に便利である。\par

\subsubsection*{・インポートプレフィックスにはsnake\_caseを使う}
\ \par

\subsubsection*{・他の識別子（変数名や関数名）にはlowerCamelCaseを使う}
クラスメンバー、トップレベル定義、変数、パラメータ、
名前付きパラメータは、
最初の単語を除く各単語の最初の文字を大文字にし、
セパレータを使用しない。\par
\clearpage

\subsubsection*{・定数には、lowerCamelCaseを使う方が良い}
新しいコードでは、enum値を含む定数変数にlowerCamelCaseを使用する。\par

次の場合のように、
既存のコードとの一貫性のためにSCREAMING\_CAPSを使用できる。\par
\begin{itemize}
  \item[$\blacktriangleright$]SCREAMING\_CAPS既に使用しているファイル
        またはライブラリにコードを追加する場合\par
  \item[$\blacktriangleright$]Javaコードと並行するDartコードを生成する場合、
        例えば、protobufsから生成された列挙されたタイプで
\end{itemize}



\subsubsection*{・3文字以上の頭文字と略語は1文字目を大文字にする（普通の単語のように扱う）}
\ \par



\subsubsection*{・未使用のコールバック引数には\ 、や\ \_\ を使用した方が良い}
\begin{itemize}
  \item[-] 関数に複数の未使用のパラメータがある場合は、
        名前の衝突を避けるために追加のアンダースコア（\_）を使用する。
\end{itemize}
\ \par


\subsubsection*{・プライベートではない識別子はアンダースコア（\_）から始めない}
\ \par
\subsubsection*{・インデント}
\begin{itemize}
  \item[-] 2スペース\par

  \item[-] 継続インデントは、4スペース\par


\end{itemize}
\ \par

\subsubsection*{・空行のインデント}
\ \par

\subsubsection*{・改行}
\begin{itemize}
  \item[-] 空行の場合はインデントを維持しない（スペースを削除する）\par

  \item[-] 引数が2つ以上のWidgetは複数行で書く\par

  \item[-] Widgetの引数が1つでも、その子Widgetが複数行の場合は、改行。
        （1行でWidgetを書くのは、Widgetツリーの末端のみ）\par

  \item[-] 例外として、Edgelnsetsは複数の引数があっても１行で良い\par

\end{itemize}
\ \par

\subsubsection*{・括弧の位置}
\begin{itemize}
  \item[-] 括弧の終了は、括弧の開始の後に改行しない場合は、同じ行で閉じる。\par


  \item[-] 括弧の開始の後に改行した場合、
        中身の最終行の次の行で括弧を閉じる。
        尚、インデントはブロック開始時と同じにする。\par


  \item[-] 関数の呼び出しで「名前なし引数」で開始する場合、
        同じ行でWidgetの括弧を開始する。
        Widgetの閉じる括弧と関数呼び出しの閉じる括弧は同じ行。\par

  \item[-] 関数の呼び出しで「名前あり引数」の場合は、
        Widgetと同様に記述。\par

\end{itemize}
\ \par
\clearpage

\subsubsection*{・コメント}
\begin{itemize}
  \item[$\blacktriangleright$] \verb|//|\ を使用
  \item[$\blacktriangleright$] コメント開始\ \verb|//|\ とコメント本文のは半角スペースひとつ開ける
\end{itemize}\par
\ \par


\subsubsection*{・クラス}
\begin{itemize}
  \item[-] new \verb|/| const\par
        オブジェクト生成時の new は省略する。 const は省略しない。\par


  \item[-] var \verb|/| final\par
        初期化値の代入以後で変更（再代入）しない変数は、
        final で宣言する。\par

\end{itemize}
\ \par
\subsubsection*{・三項演算子}
\begin{itemize}
  \item[-] 2つの値のどちらかを得るだけの場合など、
        単純な場合は使用して良い。
        結果にさらに式を含まないようにする。\par

  \item[-] 三項演算子のなかで使うものを予め計算するか、
        別にメソッドなどに切り分ける。\par

\end{itemize}



\ \par



\subsection*{順序}
\subsubsection*{・dart：から始まるインポートは他のインポートよりも前に置く}



\subsubsection*{・package：から始まるインポートは相対的なインポートよりも前に置く}


\subsubsection*{・エクスポートはインポートとは別のセクションで指定する＞}


\subsubsection*{・インポートやエクスポートは辞書順に並べ替える＞}
\ \par



\subsection*{整形}
\subsubsection*{・1行が81文字以上になることを避ける}

\subsubsection*{・ifやforの\{\}を省略しない}


\verb|-|例外として、ifにelseがなく、
かつ、全てが1行に収まる場合は{}は省略しても良い。\par
\ \par

\subsection*{・ツール}
コード作成時に自動でこれらのDart規約に従うように、
以下のツールを導入し調整を行う。\par
\begin{itemize}
  \item[-] VScodeの設定である”format on save”を有効化にし、自動的にフォーマットを行う
  \item[-] dart fix \verb|--|applyでコーディング規約に沿うように変換を行うコマンド
  \item[-] flutter analyzeでコーディング規約違反の有無を表示するコマンド
\end{itemize}

\clearpage
